\documentclass[10pt,a4paper]{amsart}
\usepackage{amsmath,amssymb,amsthm,graphicx}
% Retained minimal preamble for the Stacks Project chapter more-algebra.tex (More on Algebra).
% The chapter file pinned-original-source.tex carries only \input{preamble} on its first line, so no
% preamble text lives inside the chapter itself.  The Stacks Project uses one single shared preamble for
% every chapter; the snapshot stages it identically as sources/stacks-chow/preamble.tex and
% sources/stacks-etale-cohomology/preamble.tex.  The block below is copied mechanically, line for line,
% from sources/stacks-chow/preamble.tex lines 14-18 (the Xy-pic setup) and 154-217 (the native theorem
% declarations together with the mathematical macros).  Nothing was added, removed or redefined; the
% wrapper class and the amsmath/amssymb/amsthm/graphicx base packages are supplied by the assembler.
% Chapter: more-algebra.tex -- More on Algebra.
% For commutative diagrams we use Xy-pic
\usepackage[all]{xy}

% We use 2cell for 2-commutative diagrams.
\xyoption{2cell}

% Theorem environments.
%
\theoremstyle{plain}
\newtheorem{theorem}[subsection]{Theorem}
\newtheorem{proposition}[subsection]{Proposition}
\newtheorem{lemma}[subsection]{Lemma}

\theoremstyle{definition}
\newtheorem{definition}[subsection]{Definition}
\newtheorem{example}[subsection]{Example}
\newtheorem{exercise}[subsection]{Exercise}
\newtheorem{situation}[subsection]{Situation}

\theoremstyle{remark}
\newtheorem{remark}[subsection]{Remark}
\newtheorem{remarks}[subsection]{Remarks}

\numberwithin{equation}{subsection}

% Macros
%
\def\lim{\mathop{\mathrm{lim}}\nolimits}
\def\colim{\mathop{\mathrm{colim}}\nolimits}
\def\Spec{\mathop{\mathrm{Spec}}}
\def\Hom{\mathop{\mathrm{Hom}}\nolimits}
\def\Ext{\mathop{\mathrm{Ext}}\nolimits}
\def\SheafHom{\mathop{\mathcal{H}\!\mathit{om}}\nolimits}
\def\SheafExt{\mathop{\mathcal{E}\!\mathit{xt}}\nolimits}
\def\Sch{\mathit{Sch}}
\def\Mor{\mathop{\mathrm{Mor}}\nolimits}
\def\Ob{\mathop{\mathrm{Ob}}\nolimits}
\def\Sh{\mathop{\mathit{Sh}}\nolimits}
\def\NL{\mathop{N\!L}\nolimits}
\def\CH{\mathop{\mathrm{CH}}\nolimits}
\def\proetale{{pro\text{-}\acute{e}tale}}
\def\etale{{\acute{e}tale}}
\def\QCoh{\mathit{QCoh}}
\def\Ker{\mathop{\mathrm{Ker}}}
\def\Im{\mathop{\mathrm{Im}}}
\def\Coker{\mathop{\mathrm{Coker}}}
\def\Coim{\mathop{\mathrm{Coim}}}

% Boxtimes
%
\DeclareMathSymbol{\boxtimes}{\mathbin}{AMSa}{"02}

%
% Macros for moduli stacks/spaces
%
\def\QCohstack{\mathcal{QC}\!\mathit{oh}}
\def\Cohstack{\mathcal{C}\!\mathit{oh}}
\def\Spacesstack{\mathcal{S}\!\mathit{paces}}
\def\Quotfunctor{\mathrm{Quot}}
\def\Hilbfunctor{\mathrm{Hilb}}
\def\Curvesstack{\mathcal{C}\!\mathit{urves}}
\def\Polarizedstack{\mathcal{P}\!\mathit{olarized}}
\def\Complexesstack{\mathcal{C}\!\mathit{omplexes}}
% \Pic is the operator that assigns to X its picard group, usage \Pic(X)
% \Picardstack_{X/B} denotes the Picard stack of X over B
% \Picardfunctor_{X/B} denotes the Picard functor of X over B
\def\Pic{\mathop{\mathrm{Pic}}\nolimits}
\def\Picardstack{\mathcal{P}\!\mathit{ic}}
\def\Picardfunctor{\mathrm{Pic}}
\def\Deformationcategory{\mathcal{D}\!\mathit{ef}}

% --- appended 2026-10-12 (Chief): the project own definitions of the reference/slogan/history
% display environments, copied mechanically from the shared Stacks preamble lines 8-12
% (\usepackage{verbatim} plus the three \newenvironment lines, which the chapter preamble keeps
% verbatim from the project preamble). Nothing added, removed or redefined: these environments
% carry no mathematical content and the project itself typesets them as comments. Without these
% three lines a counted statement whose own text uses them fails with Environment slogan undefined.
\usepackage{verbatim}
\newenvironment{reference}{\comment}{\endcomment}
\newenvironment{slogan}{\comment}{\endcomment}
\newenvironment{history}{\comment}{\endcomment}
\usepackage{cleveref}
\crefname{originallemma}{original lemma}{original lemmas}
\Crefname{originallemma}{Original lemma}{Original lemmas}
\makeatletter
\newlabel{more-algebra-lemma-limit-henselian}{{0CT7}{source}{}{sourcepointer-1}{}}
\newlabel{more-algebra-lemma-limit-henselian@cref}{{[originallemma][0][]0CT7}{[1][0][]source}{}{}{}}
\makeatother
\crefname{originallemma}{original lemma}{original lemmas}
\Crefname{originallemma}{Original lemma}{Original lemmas}
\makeatletter
\newlabel{algebra-lemma-locally-nilpotent-unit}{{0AMG}{source}{}{sourcepointer-2}{}}
\newlabel{algebra-lemma-locally-nilpotent-unit@cref}{{[originallemma][0][]0AMG}{[1][0][]source}{}{}{}}
\makeatother
\crefname{originallemma}{original lemma}{original lemmas}
\Crefname{originallemma}{Original lemma}{Original lemmas}
\makeatletter
\newlabel{algebra-lemma-contained-in-radical}{{0AME}{source}{}{sourcepointer-3}{}}
\newlabel{algebra-lemma-contained-in-radical@cref}{{[originallemma][0][]0AME}{[1][0][]source}{}{}{}}
\makeatother
\crefname{originallemma}{original lemma}{original lemmas}
\Crefname{originallemma}{Original lemma}{Original lemmas}
\makeatletter
\newlabel{lemma-locally-nilpotent-henselian}{{0ALI}{source}{}{sourcepointer-4}{}}
\newlabel{lemma-locally-nilpotent-henselian@cref}{{[originallemma][0][]0ALI}{[1][0][]source}{}{}{}}
\makeatother
\title{Inverse limits of rings along surjections with locally nilpotent kernels are henselian pairs}
\author{The Stacks Project Authors}
\date{stacks-project:more-algebra.tex:local-snapshot-20261012; GNU FDL 1.2 or later}
\begin{document}
\maketitle
\setcounter{subsection}{1}
\noindent\textit{Editorial note.} All delivered mathematical text is the author's own, copied from the pinned Stacks Project chapter file more-algebra.tex; no statement or proof step was rewritten, reordered or removed.\par\medskip
\noindent\textit{Editorial note.} The chapter file pinned-original-source.tex carries only \textbackslash{}input\{preamble\} on its first line, so no preamble text lives inside the chapter. The retained file sources/stacks-more-algebra/retained-preamble.tex is a mechanical, line-for-line copy of the project's single shared preamble as staged in the snapshot (Xy-pic setup plus theorem declarations and mathematical macros). Nothing was added, removed or redefined, and the excerpt serves only as preamble, never as a span coordinate.\par\medskip
\noindent\textit{Editorial note.} Both counted spans start at the \textbackslash{}begin line of their environment, so the statement and the author's proof are delivered whole; surrounding prose of the section is deliberately outside the counted spans and is carried as standard prerequisites instead.\par\medskip
\noindent\textit{Editorial note.} The citation fields (source\_id, source\_version, authors, source\_url, license.name, source\_title) are identical to those of every other card this lane has produced from this chapter, so the batch selector treats them as one source.\par\medskip
\noindent\textit{Allowed standard prerequisites (editorial).}\begin{itemize}
\item An element that is a unit modulo a locally nilpotent kernel is a unit. Algebra, Lemma \textbackslash{}ref\{algebra-lemma-locally-nilpotent-unit\}, cited at proof lines 2460-2461. Shows that elements which are units at some stage are units in the limit.
\item Elements of one plus an ideal consisting of such units force the ideal into the Jacobson radical. Algebra, Lemma \textbackslash{}ref\{algebra-lemma-contained-in-radical\}, cited at proof lines 2465-2466. Puts the kernel ideal inside the Jacobson radical, one half of the henselian condition.
\item Lifting of factorisations of monic polynomials through surjections with locally nilpotent kernels. Lemma \textbackslash{}ref\{lemma-locally-nilpotent-henselian\}, cited at proof lines 2471-2472. Provides the successive lifting of the factorisation of the reduction of the monic polynomial.
\end{itemize}\medskip
\ifdefined\hypertarget\hypertarget{sourcepointer-1}{}\fi
\noindent\textit{Source reference.} Original lemma 0CT7: focused-seeds/AGP/sources/stacks-snapshot-20261011/tags.txt, source lines 16057--16057. The lemma's own canonical tag, quoted so that readers can resolve the label inside the delivered text; it points at the counted statement itself and carries no proof content.\par\medskip
\ifdefined\hypertarget\hypertarget{sourcepointer-2}{}\fi
\noindent\textit{Source reference.} Original lemma 0AMG: focused-seeds/AGP/sources/stacks-snapshot-20261011/tags.txt, source lines 13404--13404. Navigation pointer to a statement the counted proof refers to but does not contain; nothing load-bearing is delegated to it.\par\medskip
\ifdefined\hypertarget\hypertarget{sourcepointer-3}{}\fi
\noindent\textit{Source reference.} Original lemma 0AME: focused-seeds/AGP/sources/stacks-snapshot-20261011/tags.txt, source lines 13402--13402. Navigation pointer to a statement the counted proof refers to but does not contain; nothing load-bearing is delegated to it.\par\medskip
\ifdefined\hypertarget\hypertarget{sourcepointer-4}{}\fi
\noindent\textit{Source reference.} Original lemma 0ALI: focused-seeds/AGP/sources/stacks-snapshot-20261011/tags.txt, source lines 13371--13371. Navigation pointer to a statement the counted proof refers to but does not contain; nothing load-bearing is delegated to it.\par\medskip
\begin{lemma}
\label{lemma-limit-henselian}
Let $A = \lim A_n$ where $(A_n)$ is an inverse system of rings
whose transition maps are surjective and have locally nilpotent kernels.
Then $(A, I_n)$ is a henselian pair, where $I_n = \Ker(A \to A_n)$.
\end{lemma}

\begin{proof}
Fix $n$. Let $a \in A$ be an element which maps to $1$ in $A_n$.
By Algebra, Lemma \ref{algebra-lemma-locally-nilpotent-unit}
we see that $a$ maps to a unit in $A_m$ for all $m \geq n$.
Hence $a$ is a unit in $A$. Thus by
Algebra, Lemma \ref{algebra-lemma-contained-in-radical}
the ideal $I_n$ is contained in the Jacobson radical of $A$.
Let $f \in A[T]$ be a monic polynomial and let
$\overline{f} = g_nh_n$ be a factorization
of $\overline{f} = f \bmod I_n$ with $g_n, h_n \in A_n[T]$ monic
generating the unit ideal in $A_n[T]$. By
Lemma \ref{lemma-locally-nilpotent-henselian}
we can successively lift this factorization to
$f \bmod I_m = g_m h_m$ with $g_m, h_m$ monic
in $A_m[T]$ for all $m \geq n$. At each step we have
to verify that our lifts $g_m, h_m$ generate the unit ideal in
$A_n[T]$; this follows from the corresponding fact for
$g_n, h_n$ and the fact that $\Spec(A_n[T]) = \Spec(A_m[T])$
because the kernel of $A_m \to A_n$ is locally nilpotent.
As $A = \lim A_m$ this finishes the proof.
\end{proof}

\end{document}
